\section{Introduction}
\label{sec:intro}

The modern machine learning paradigm heavily relies on pre-training large-scale foundational models and subsequently fine-tuning them on specialized downstream tasks. While this approach has driven remarkable progress, deploying a separate, full-scale fine-tuned expert network for every target task becomes prohibitively expensive under tight deployment constraints, such as on-device and edge-computing applications. Maintaining multiple separate models leads to linear scaling of memory footprints and storage overhead, which is unsustainable.

To address this challenge, \emph{model merging} has emerged as an active and promising research direction. Instead of keeping multiple expert checkpoints or retraining a joint multi-task model from scratch (which is computationally expensive and requires access to all joint datasets), model merging aims to consolidate multiple task-specialized models directly into a single unified backbone. A standard baseline in this area is \emph{Task Arithmetic} \cite{ilharco2022editing}, where task-specific parameter shifts (called task vectors, $V_k = W_k - W_{base}$) are linearly scaled and added directly to the pre-trained base model weights.

Despite its simplicity and utility, static model merging is fundamentally limited. Static blending techniques (such as uniform task arithmetic) apply the exact same task weights across all inputs and layers. This assumption is structurally restrictive: task-specific parameters often contain highly conflicting representations. Linearly blending them without regard to layer-wise or sample-wise specificities leads to severe representational interference, parameter collision, and dramatic accuracy collapse on individual tasks.

To mitigate this interference, recent works have proposed dynamic merging approaches. These include input-dependent linear routers that compute layer-wise coefficients dynamically, or wavefunction superposition frameworks (such as QWS-Merge) that project inputs to a phase space. However, these methods either employ unconstrained, over-parameterized routing functions that overfit to small calibration sets (the Overfitting-Optimizer Paradox), or rely on flat, linear feature projections that fail to capture the highly non-linear, multi-scale representational boundaries of deep neural networks.

In this work, we present a radical departure from traditional flat parameter blending. Guided by a visionary research philosophy, we ask: \emph{What if the sequence of layers in a deep neural network is not merely a feed-forward pipeline, but the temporal evolution of a chaotic dynamical system?} 

We introduce \textbf{ChaosMerge} (Chaos-Theoretic Attractor Merging), a novel and highly effective framework for parameter-space model fusion. Instead of using classical feed-forward neural layers to compute merging weights, ChaosMerge frames the layer-wise model-merging coefficients as trajectories of a chaotic Coupled Map Lattice (CML) driven by a Logistic Map. Our key insight is that treating layer depth $l$ as discrete time steps in a CML allows merging coefficients to self-organize dynamically. Input representations are mapped to a low-dimensional unit sphere and injected as logit-space perturbations to steer the chaotic orbit. 

However, propagating fully chaotic Logistic Maps recursively over 14 layers in a deep recurrent lattice leads to an exponential gradient multiplier, resulting in severe gradient explosion up to $4^{14} \approx 2.68 \times 10^8$. To resolve this fundamental bottleneck, we propose a \textbf{Gated Coupled Map Lattice (G-CML)}, which incorporates a learned layer-wise gating coefficient acting as a residual skip-connection. G-CML successfully tames gradient explosion, providing a stable additive pathway for gradients to propagate back cleanly, enabling robust optimization with standard first-order optimizers. Furthermore, we address the batch-averaging contradiction of dynamic merging by introducing \textbf{Task-Specific Dynamic Routing} which utilizes task-level centroid features at test-time. This extracts task-specific coefficients directly, preventing individual task signatures from being washed out in heterogeneous batches while ensuring negligible test-time computational and memory-swapping latency.

Empirical evaluations on visual classification tasks (MNIST, FashionMNIST, CIFAR-10, SVHN) using a Vision Transformer backbone demonstrate that ChaosMerge achieves highly competitive results. With only 384 parameters, ChaosMerge significantly outperforms competitive static merging baselines (such as Task Arithmetic and AdaMerging) and achieves performance close to static supervised tuning and unconstrained linear routers with nearly $30\times$ more parameters, validating our hypothesis that stabilized chaotic dynamics provide exceptional structural regularization and dynamic separation.

\subsection{Our Contributions}
Our primary contributions are summarized as follows:
\begin{itemize}
    \item \textbf{A Gated Chaotic Merging Paradigm (G-CML):} We propose the first framework combining chaos theory and model merging, introducing the Gated Coupled Map Lattice (G-CML) to successfully tame chaotic gradient explosion and stabilize parameter-space optimization.
    \item \textbf{Task-Specific Dynamic Routing:} We resolve the batch-averaging contradiction of dynamic merging by performing task-specific coefficient routing, preventing individual task characteristics from washing out and enabling true dynamic input-adaptation.
    \item \textbf{Extremely Compact Parameter Footprint:} We achieve highly parameter-efficient dynamic model merging using a physically regularized lattice of only 384 parameters, successfully avoiding the Overfitting-Optimizer Paradox.
    \item \textbf{Outstanding Empirical Results:} We present a rigorous comparison of ChaosMerge against five competitive static and dynamic baselines on multi-task visual classification, showing that G-CML boosts average merging accuracy by \textbf{+18.60\%} absolute over the original chaotic baseline while achieving highly competitive performance compared to over-parameterized dynamic baselines with a fraction of their parameter footprints.
\end{itemize}
